% English translation of questions/5778/semA-moedA/q1b.tex (metadata is taken from the Hebrew file).

\begin{question}
Compute the following limit:
the limit of the function $\displaystyle \lim_{x\to\pi/2}(\tan x)\cdot\arctan\left(x-\frac{\pi}{2}\right)$.
\end{question}

\begin{hint}
Write $\tan x=\frac{1}{\cot x}$ and obtain an expression of the form $\frac00$.
\end{hint}

\begin{solution}
For $x$ close to $\frac\pi2$, $x\neq\frac\pi2$, we have $\tan x=\frac{1}{\cot x}$, hence
\[ (\tan x)\cdot\arctan\left(x-\frac{\pi}{2}\right)=\frac{\arctan\left(x-\frac{\pi}{2}\right)}{\cot x}. \]
As $x\to\frac\pi2$: the numerator tends to $\arctan 0=0$ and the denominator tends to $\cot\frac\pi2=0$ (both functions are continuous). This is an expression of the form $\frac00$, both functions are differentiable in a punctured neighborhood of $\frac\pi2$, and the derivative of the denominator $-\frac{1}{\sin^2x}\neq 0$ there, so L'Hôpital's rule can be applied:
\[ \lim_{x\to\pi/2}\frac{\arctan\left(x-\frac{\pi}{2}\right)}{\cot x}
   =\lim_{x\to\pi/2}\frac{\dfrac{1}{1+\left(x-\frac\pi2\right)^2}}{-\dfrac{1}{\sin^2x}}
   =\lim_{x\to\pi/2}\left(-\frac{\sin^2x}{1+\left(x-\frac\pi2\right)^2}\right)=-\frac{1}{1+0}=-1, \]
where the last limit is computed by continuity ($\sin^2\frac\pi2=1$). Since the limit of the quotient of the derivatives exists, by L'Hôpital's rule
\[ \lim_{x\to\pi/2}(\tan x)\cdot\arctan\left(x-\frac{\pi}{2}\right)=-1. \]
(Another way: $\arctan t\sim t$ as $t\to0$, and $(x-\frac\pi2)\tan x=-\frac{(x-\frac\pi2)\cos(x-\frac\pi2)}{\sin(x-\frac\pi2)}\to-1$.)
\end{solution}
